\section{The Arrow of Time Is Not in the
Ensemble}\label{the-arrow-of-time-is-not-in-the-ensemble}

\subsection{An Audited Toy-Model Quest across Equilibrium Matrices,
Sequential Growth, and Causal Sets (SGOED
V7)}\label{an-audited-toy-model-quest-across-equilibrium-matrices-sequential-growth-and-causal-sets-sgoed-v7}

\textbf{Author:} Sutipong Chanpengpad \emph{Independent Researcher,
Chiang Rai, Thailand} Email: \texttt{dhammawatthumpra@gmail.com} ORCID:
\href{https://orcid.org/0009-0001-4069-8576}{0009-0001-4069-8576}

\textbf{Date:} September 2026

\begin{center}\rule{0.5\linewidth}{0.5pt}\end{center}

\subsection{Abstract}\label{abstract}

We ask a focused question in a toy universe of coupled
matrix/graph/causal-set models: can the arrow of time arise from the
\emph{final equilibrium state} of a relational structure, however it is
coupled? Using a disciplined audit (reproduce, significance, null test,
labeling test, thermalization check, mechanism --- seed counts stated
per result, \(\ge 10\) wherever per-seed scatter is non-trivial), we
demonstrate that it cannot: \textbf{eleven structural observables}
(asymmetry ratio, sign counts, Myrheim--Meyer dimension, spectral
dimension, cycle ratio, time-reversal, order statistics, eigenvalue
repulsion, entropy rate, \ldots) are all time-symmetric or
null-compatible at equilibrium Monte Carlo. In contrast, three
\emph{process-based} constructions produce robust directional time:

\begin{enumerate}
\def\labelenumi{\arabic{enumi}.}
\tightlist
\item
  \textbf{Past Hypothesis} --- a low-entropy (rank-1) initial state
  relaxes with \(dS/S_0 = 0.98\), roughly \(100\times\) the relaxation
  of a random initial state.
\item
  \textbf{Sequential growth} (Rideout--Sorkin style) --- birth order
  with frozen past yields chain inheritance \(1.000\pm0.000\),
  deterministic and scale-free (\(M=16,32\)); the mechanism is a
  contraction map induced by alignment to the past-mean direction
  (prefix-sum coupling).
\item
  \textbf{Non-equilibrium upwind transport} --- a directed source→sink
  flow produces a \(5000\times\)--\(6000\times\) extent decay with
  current \(J\approx +5.77\pm0.02\) (8 seeds), reversal-symmetric
  (\(-6.03\)) and null-discriminating (\(J\approx 0\)).
\end{enumerate}

We then audit the ``emergent 4D'' question that drove the project's
excitement into a dead end. Euclidean IKKT-like actions at \(D=4,10\)
are perfectly isotropic without an engineered steering vector
\(\hat v\); a pseudo-Euclidean signature produces a stable \(1{+}3\)
split only as an engineered \(\eta\) (spatial isotropy \(10^{-3}\));
Lorentzian real-time bosonic dynamics is unbounded below (the regulator
in the literature is the fermion determinant); random-percolation causal
sets give ``\(d\approx4\)'' only as a \textbf{finite-size artifact of
\(N=250\)} that fails a scale audit to \(N=1000\) under both fixed-\(p\)
and fixed-\(k\) scalings; and spectral dimension measured on flow
networks tracks drift/transience rather than geometry. The one
structural discovery is that a \textbf{relational light-cone law}
(matrix-distance \(\le c\,\Delta t\) plus direction alignment) makes the
causal dimension \emph{scale-invariant} (\(\mathrm{std}\) across \(N\)
of \(0.002\)--\(0.01\), versus \(0.26\)--\(0.82\) for percolation) --- a
necessary condition for a dimensionful causal set --- although the
absolute dimension remains parameter-controlled (diffusion/light-speed
aspect), i.e., engineered.

\textbf{Bottom line:} at toy scale, the arrow of time lives in the
\emph{process} (initial condition or asymmetric construction), not in
the equilibrium ensemble; and 4D does not emerge spontaneously in any of
the audited frameworks. The audit itself --- a six-step verification
protocol and a trap catalog --- is the project's transferable
contribution, shipped as a runnable toolbox
(\texttt{code/audit\_gates.py}, Supplement A) with the full v8--v14
evidence slice (\texttt{audit\_evidence/}).

\textbf{Keywords:} arrow of time, matrix models, sequential growth,
causal sets, spectral dimension, Monte Carlo audit, past hypothesis

\textbf{Important Caveats:} toy-level demonstration with small sizes
(\(N\le 1000\) events, matrices \(N\le 8\)). ``Real time'' here means
\emph{robust directional asymmetry in a specified observable},
reproduced deterministically across seeds --- not a claim about
fundamental physics.

\begin{center}\rule{0.5\linewidth}{0.5pt}\end{center}

\subsection{Table of Contents}\label{table-of-contents}

\begin{itemize}
\tightlist
\item
  \hyperref[1-introduction]{1. Introduction}
\item
  \hyperref[2-models-and-observables]{2. Models and Observables}

  \begin{itemize}
  \tightlist
  \item
    \hyperref[21-the-matrix-unit-v7-core]{2.1 The matrix unit (v7 core)}
  \item
    \hyperref[22-growth-transport-and-posets]{2.2 Growth, transport, and
    posets}
  \item
    \hyperref[23-observables]{2.3 Observables}
  \end{itemize}
\item
  \hyperref[3-the-six-step-audit]{3. The Six-Step Audit}
\item
  \hyperref[4-results]{4. Results}

  \begin{itemize}
  \tightlist
  \item
    \hyperref[41-equilibrium-is-time-symmetric]{4.1 Equilibrium is
    time-symmetric}
  \item
    \hyperref[42-route-1-past-hypothesis]{4.2 Route 1: Past Hypothesis}
  \item
    \hyperref[43-route-2-sequential-growth-and-its-mechanism]{4.3 Route
    2: Sequential growth and its mechanism}
  \item
    \hyperref[44-route-3-non-equilibrium-upwind-transport]{4.4 Route 3:
    Non-equilibrium upwind transport}
  \item
    \hyperref[45-the-4d-question-matrix-spontaneous-symmetry-breaking]{4.5
    The 4D question: matrix spontaneous symmetry breaking}
  \item
    \hyperref[46-causal-set-dimensions-percolation-negative-light-cone-invariance]{4.6
    Causal-set dimensions: percolation negative, light-cone invariance}
  \item
    \hyperref[47-spectral-dimension-on-flow-networks]{4.7 Spectral
    dimension on flow networks}
  \end{itemize}
\item
  \hyperref[5-discussion]{5. Discussion}
\item
  \hyperref[6-conclusion]{6. Conclusion}
\item
  \hyperref[appendix-a-reproducibility]{Appendix A: Reproducibility}
\item
  \hyperref[appendix-b-trap-catalog]{Appendix B: Trap Catalog}
\item
  \hyperref[references]{References}
\end{itemize}

\begin{center}\rule{0.5\linewidth}{0.5pt}\end{center}

\subsection{1. Introduction}\label{introduction}

The question of where the arrow of time comes from has been posed across
philosophy {[}\hyperref[ref-1]{1}, \hyperref[ref-9]{9},
\hyperref[ref-10]{10}{]}, thermodynamics, and relational physics
{[}\hyperref[ref-2]{2}, \hyperref[ref-5]{5}{]}. Our working hypothesis,
following the Structural Time Framework {[}\hyperref[ref-6]{6}{]} and
its SGOED implementation {[}\hyperref[ref-7]{7}{]}, is relational: time
might emerge from the \emph{structure} of a pre-temporal system rather
than being inserted as a parameter.

The project evolved through many constructions --- matrix models with an
observer system {[}\hyperref[ref-7]{7}{]}, graphs, hypergraphs,
ecosystems, hybrids --- and each produced metrics that ``looked
directional.'' Exactly eleven of them collapsed under scrutiny: they
were histogram artifacts, labeling-dependent, null-compatible, or
thermalization-limited. Two conclusions emerged as robust:

\begin{enumerate}
\def\labelenumi{\arabic{enumi}.}
\tightlist
\item
  \textbf{Equilibrium Monte Carlo is time-symmetric.} No coupling of an
  equilibrium ensemble produces an arrow from the final state.
\item
  \textbf{The arrow requires process:} either a special initial
  condition (past hypothesis) or a construction whose updates are
  asymmetric by design (birth order, directed transport).
\end{enumerate}

This manuscript reports the audited results of that quest, plus the
systematic audit of the ``emergent 4D'' question that followed, and the
one structural discovery it produced: relational light-cone locality as
the condition for \emph{scale-invariant} causal dimension.

\begin{quote}
\textbf{Important Disclaimer:} A toy-level computational exploration
with small sizes. ``Time'' below means \emph{robust, reproducible
directional asymmetry of an observable along a constructed axis}, not a
claim about fundamental physics. Everything reported is reproduced from
public scripts (Appendix \hyperref[appendix-a-reproducibility]{A}).
\end{quote}

\begin{center}\rule{0.5\linewidth}{0.5pt}\end{center}

\subsection{2. Models and Observables}\label{models-and-observables}

\subsubsection{2.1 The matrix unit (v7
core)}\label{the-matrix-unit-v7-core}

A unit holds real symmetric matrices \(X = (X^1,\dots,X^D)\) of size
\(N\times N\) (``system'') and \(Y = (Y^1,\dots,Y^d)\) (``observer'').
The local action \(S_X + S_Y + S_{\mathrm{coup}}\) (all terms of the
IKKT-like family {[}\hyperref[ref-15]{15}{]} with stability gates) is

\[S = \sum_{\mu<\nu}\mathrm{Tr}([X_\mu,X_\nu][X_\mu,X_\nu]^T) + \sum_\mu \lambda\big(\mathrm{Tr}(X_\mu^2)-Nr_0^2\big)^2 \ -\ g_{XY}\sum_\mu \hat v_\mu^2\,\mathrm{Tr}(X_\mu^4),\]

where
\(\hat v = \mathrm{normalize}(\mathrm{Tr}(Y^1),\dots,\mathrm{Tr}(Y^d))\)
is the observer direction. The coupling
\(-\hat v_\mu^2\mathrm{Tr}(X_\mu^4)\) drives exactly the dimension(s)
aligned with \(\hat v\) to large extent \(\mathrm{Tr}(X_\mu^2)/N\) --- a
\emph{discrete} winner-take-all choice among \(D\) dimensions
(condensation ratio \(\lambda_{\max}/\lambda_{2\mathrm{nd}}\approx 23\),
\(\sim 60\sigma\) over baseline {[}\hyperref[ref-7]{7}{]}), with
first-order bistability/hysteresis (audited scripts
\texttt{AUDIT\_v7\_hysteresis*} in \texttt{audit\_evidence/}; see also
the v6 preprint). This mechanism is the reproducible core of the
project; it picks \emph{a} direction but does not create \emph{flow}.

\subsubsection{2.2 Growth, transport, and
posets}\label{growth-transport-and-posets}

Three process constructions were added on top of the unit:

\begin{itemize}
\tightlist
\item
  \textbf{Sequential growth (CSG-style {[}\hyperref[ref-11]{11}{]}):}
  units are born one at a time; past units are frozen; the newborn
  thermalizes against the frozen past with coupling
  \(S_{\mathrm{inter}} = -g\,\hat v_u\cdot \big(\sum_{j<u}\hat v_j\big)\)
  (prefix-sum --- an exact \(O(1)\) reformulation of pairwise
  alignment).
\item
  \textbf{Upwind transport (Langevin, non-MC):} \(M\) units on a chain;
  a drive pumps unit \(0\), a sink absorbs at unit \(M-1\), and the
  coupling transfers extent only forward
  (\(F_u \propto g_t X_{u-1} - g_t X_u\)); updates are overdamped
  Langevin with temperature \(T\).
\item
  \textbf{Causal sets / posets} {[}\hyperref[ref-12]{12}{]}: relations
  between births are defined either probabilistically (Bernoulli links +
  transitive closure), or geometrically from matrix state --- unit
  \(u\prec k\) iff \(|\hat v_u\cdot \hat v_k|>\theta\) \textbf{and}
  \(D_{\mathrm{space}}(u,k)\le (c\,\Delta t)^2\) with
  \(D_{\mathrm{space}} = \frac{1}{N}\mathrm{Tr}((X_u-X_k)^2)\)
  (\textbf{relational light-cone law}).
\end{itemize}

\subsubsection{2.3 Observables}\label{observables}

Directionality metrics: extent ratio
\(R=\mathrm{ext}(X_{\max})/\overline{\mathrm{ext}}(X_{\neq\max})\);
sign-count \(D\); relational dimension via relation fraction
\(\rho = C/\binom{N}{2}\) calibrated against synthetic Minkowski
sprinklings (\(d=1..8\), per-\(N\) recalibration to cancel finite-size
box effects; estimator after Myrheim {[}\hyperref[ref-13]{13}{]});
spectral dimension \(d_s = -2\,d\ln P/d\ln t\) from the return
probability \(P(t)=\mathrm{Tr}(T^t)/N\) of the (lazy) random walk;
current
\(J = \overline{\mathrm{ext}}_{u} - \overline{\mathrm{ext}}_{u+1}\);
entropies and counts thereof.

\begin{center}\rule{0.5\linewidth}{0.5pt}\end{center}

\subsection{3. The Six-Step Audit}\label{the-six-step-audit}

Every observable in this project is passed through six gates before
being reported. The gates, with their pass criteria (implemented as a
runnable toolbox, \texttt{code/audit\_gates.py}, Supplement A):

{\def\LTcaptype{none} % do not increment counter
\begin{longtable}[]{@{}
  >{\raggedright\arraybackslash}p{(\linewidth - 2\tabcolsep) * \real{0.5000}}
  >{\raggedright\arraybackslash}p{(\linewidth - 2\tabcolsep) * \real{0.5000}}@{}}
\toprule\noalign{}
\begin{minipage}[b]{\linewidth}\raggedright
gate
\end{minipage} & \begin{minipage}[b]{\linewidth}\raggedright
pass criterion
\end{minipage} \\
\midrule\noalign{}
\endhead
\bottomrule\noalign{}
\endlastfoot
1. Reproduce & identical numbers on re-run (same seed, to machine
precision) \\
2. Significance & \(\ge 10\) seeds; mean and std stabilize (5 seeds are
not enough --- a \(2.3\sigma\) effect fell to \(1.2\sigma\) in an
earlier branch {[}\hyperref[ref-7]{7}{]}) \\
3. Null test & the effect separates from shuffle/baseline (≳ \(3\sigma\)
over the null spread) \\
4. Labeling/permutation & the metric is invariant under relabeling of
nodes \\
5. Thermalization & the observable plateaus once \(n_{\mathrm{therm}}\)
is sufficient (this gate alone reversed two conclusions in this project
--- Section
\hyperref[43-route-2-sequential-growth-and-its-mechanism]{4.3},
\hyperref[42-route-1-past-hypothesis]{4.2}) \\
6. Mechanism & a stated explanation of \emph{why}, backed by the
measured evidence \\
\end{longtable}
}

\emph{Seed-count note.} The table's ``\(\ge 10\) seeds'' criterion
applies to observables whose per-seed scatter is non-trivial (every
dimension estimate and every null/labeling control here used
\(10\)--\(12\) seeds). The three deterministic routes report fewer seeds
\emph{because their scatter is negligible by construction}: sequential
growth inherits the past mean deterministically (origin alignment
\(0.9995\pm0.0002\) std over 6 seeds, chain \(0.9991\pm0.0002\)), and
the upwind transport current is \(J=+5.773\pm0.016\) over 8 seeds
(\(\sim0.3\%\)). For such nearly-deterministic observables, additional
seeds only confirm the mean; we state the actual counts per result in
Appendix \hyperref[appendix-a-reproducibility]{A} rather than claiming a
uniform \(10\).

Appendix \hyperref[appendix-b-trap-catalog]{B} lists the concrete
failure modes this process caught.

\begin{center}\rule{0.5\linewidth}{0.5pt}\end{center}

\subsection{4. Results}\label{results}

\subsubsection{4.1 Equilibrium is
time-symmetric}\label{equilibrium-is-time-symmetric}

In equilibrium Metropolis sampling {[}\hyperref[ref-8]{8}{]}, no
final-state observable distinguishes a direction of time. Eleven
observables were attempted across matrix, graph, hypergraph, and
ecosystem variants: extent ratio \(R\), sign-count \(D\), relational
dimension \(d_{\mathrm{MM}}\), spectral dimension \(d_s\), graph
asymmetry \(G\), cycle ratio, time-reversal statistic, order statistics,
eigenvalue repulsion, dynamic order, and entropy rate. Their failures
are catalogued in Appendix \hyperref[appendix-b-trap-catalog]{B} (e.g.,
\(R\approx0.5\) is a histogram artifact; \(D\) is labeling-dependent,
ranging \([-98,+28]\) under node permutations; \(d_{\mathrm{MM}}\) and
\(d_s\) are null-compatible with random graphs). The two genuinely
asymmetric phases --- eigenvalue condensation (discrete \(\mu\) choice)
and inter-unit clock synchronization (alignment \(1.000\pm0.001\),
null-passed) --- are real but \emph{static}: they fix a direction in
space, not in time.

\subsubsection{4.2 Route 1: Past
Hypothesis}\label{route-1-past-hypothesis}

Following the Boltzmann/Albert past-hypothesis program
{[}\hyperref[ref-16]{16}{]}: a special low-entropy initial state relaxes
toward equilibrium, and the arrow is the measure of that relaxation. In
the v7 matrix unit (\(N=4\), \(g_{XY}=0\), 3 seeds):

{\def\LTcaptype{none} % do not increment counter
\begin{longtable}[]{@{}lll@{}}
\toprule\noalign{}
initial state & \(dS/S_0\) & total \(dS\) \\
\midrule\noalign{}
\endhead
\bottomrule\noalign{}
\endlastfoot
random & 0.29 & 3 \\
\textbf{rank-1 (low entropy)} & \textbf{0.98} & \textbf{337} \\
\end{longtable}
}

The rank-1 initial condition relaxes \(\sim100\times\) more entropy than
a random one, deterministically and monotonically (\(S\): \(344\to7\)).
\emph{Mechanism:} low-entropy states have more relaxation capacity; the
arrow's strength is inherited from the initial condition, not from the
dynamics.

\subsubsection{4.3 Route 2: Sequential growth and its
mechanism}\label{route-2-sequential-growth-and-its-mechanism}

\textbf{Result.} Units born one-by-one with frozen past (coupling
\(g=20\), adequate thermalization) show:

{\def\LTcaptype{none} % do not increment counter
\begin{longtable}[]{@{}lll@{}}
\toprule\noalign{}
\(M\) & chain inheritance & origin alignment \\
\midrule\noalign{}
\endhead
\bottomrule\noalign{}
\endlastfoot
8 & \(1.000 \pm 0.000\) & \(1.000\) \\
16 & \(0.998\) & \(0.998\) \\
32 & \(0.998\) & \(0.998\) \\
\end{longtable}
}

Deterministic across seeds (identical statistics). \textbf{Mechanism
(why 1.000):} the inter-coupling aligns each newborn to the \emph{mean
direction of all past units}; the past-mean is a contraction fixed
point. Measured contraction of late units to the past mean: \(0.03\) at
\(g=20\), \(0.02\) at \(g=60\). The alignment is partial at weak
coupling (chain \(0.43\) at \(g=1\)) and saturates at strong coupling
--- a clean coupling-strength phase structure. \textbf{Scale-free} from
\(M=8\) to \(M=32\).

\textbf{Thermalization gate.} This route initially failed a reproduction
at short thermalization (\(n_{\mathrm{therm}}=30\)): drift appeared and
inheritance dropped. At \(n_{\mathrm{therm}}=100\) the drift vanished
--- a thermalization artifact, not physics.

\textbf{Past hypothesis + growth.} The handoff reported a failed attempt
to seed the origin with a special state. The correct fix combines: (a) a
rank-1 origin (\(Y^1_0=2I\) so \(\hat v_0=(1,0)\)), (b) frozen origin,
(c) sufficient thermalization (\(n_{\mathrm{therm}}=120\)). With the
origin untrained the earlier failure was \emph{thermalization
insufficiency}, not origin thermalization: at \(n_{\mathrm{therm}}=60\)
origin alignment was \(0.68\pm0.45\) (one of six seeds unconverged); at
\(120\) it is \textbf{\(0.9995\pm0.0002\)} with chain
\(0.9991\pm0.0002\) (6 seeds). A random-origin control gives
\(0.90\pm0.10\) --- the special initial condition genuinely determines
the propagated direction. Route 2 thus realizes the past-hypothesis
arrow \emph{within} a growth process.

\subsubsection{4.4 Route 3: Non-equilibrium upwind
transport}\label{route-3-non-equilibrium-upwind-transport}

Overdamped Langevin on a chain of \(M=6\) units with upwind transfer,
pump at unit 0, sink at unit \(M-1\) (parameters:
\(g_{\mathrm{drive}}=5\), \(g_{\mathrm{trans}}=1.2\),
\(g_{\mathrm{sink}}=3\), \(T=0.02\)):

{\def\LTcaptype{none} % do not increment counter
\begin{longtable}[]{@{}
  >{\raggedright\arraybackslash}p{(\linewidth - 4\tabcolsep) * \real{0.3333}}
  >{\raggedright\arraybackslash}p{(\linewidth - 4\tabcolsep) * \real{0.3333}}
  >{\raggedright\arraybackslash}p{(\linewidth - 4\tabcolsep) * \real{0.3333}}@{}}
\toprule\noalign{}
\begin{minipage}[b]{\linewidth}\raggedright
setup
\end{minipage} & \begin{minipage}[b]{\linewidth}\raggedright
extent profile \(E_0..E_5\)
\end{minipage} & \begin{minipage}[b]{\linewidth}\raggedright
current \(J\)
\end{minipage} \\
\midrule\noalign{}
\endhead
\bottomrule\noalign{}
\endlastfoot
\textbf{REAL (pump@0)} &
\([28.8,\,1.72,\,0.117,\,0.0247,\,0.020,\,0.005]\) &
\textbf{\(+5.773\pm0.016\)} \\
REVERSED (pump@5) & \(\sim[0.02,\dots,30.1]\) & \(-6.025\pm0.022\) \\
NULL (no pump) & flat \(\sim0.1\) & \(+0.003\pm0.008\) \\
\end{longtable}
}

\begin{itemize}
\tightlist
\item
  \textbf{Decay \(5849\times\)} (\(=28.8/0.005\); robust audit, 8 seeds
  \(\times\) 500 steps), dramatically larger than any equilibrium
  asymmetry.
\item
  \textbf{Deterministic:} current std is \(0.016\) over 8 seeds
  (\(\sim0.3\%\)).
\item
  \textbf{Reversal-symmetric:} pumping at the far end mirrors the
  profile and flips the current sign --- the asymmetry follows the pump,
  not a latent bias.
\item
  \textbf{Null-discriminating:} without a pump the profile is flat and
  \(J\approx0\).
\item
  \textbf{Parameter family:} scanning
  \(g_{\mathrm{trans}}\times g_{\mathrm{sink}}\) over a \(3\times3\)
  grid keeps \(J\in[+5.4,+5.9]\) and decay \(1850\)--\(4000\times\),
  always monotone in the mean.
\end{itemize}

\emph{Caveat:} the mean profile is strictly monotone, but per-seed
profiles occasionally violate monotonicity at the thermal floor
(\(\sim0.02\), three orders below \(E_0\)) --- ``monotone in practice'',
with the residual explained as noise, consistently with the thermal
floor.

\subsubsection{4.5 The 4D question: matrix spontaneous symmetry
breaking}\label{the-4d-question-matrix-spontaneous-symmetry-breaking}

Motivated by the Lorentzian IKKT programs {[}\hyperref[ref-3]{3},
\hyperref[ref-4]{4}, \hyperref[ref-15]{15}{]}, we asked whether the
matrix unit itself can spontaneously pick a \(1{+}3\) split at \(D=4\)
(and \(1{+}9\) at \(D=10\)) \emph{without} the engineered steering
\(\hat v\). All four probes were audited at \(\ge 10\) seeds:

\begin{enumerate}
\def\labelenumi{\arabic{enumi}.}
\tightlist
\item
  \textbf{Equilibrium, no steering (\(g=0\) or symmetric driver):}
  perfectly isotropic at \(D=4\) (\(R=1.06\pm0.03\)) and \(D=10\)
  (\(R=1.07\pm0.02\)). The ``direction choice'' of the v7 action is
  entirely \(\hat v\)-engineered (the D=2 control reads
  \(R=4.43\pm0.16\), matching the documented value).
\item
  \textbf{Pseudo-Euclidean signature}
  (\(\eta=\mathrm{diag}(-1,+1,\dots)\), so space--space commutators
  enter with opposite sign): a stable \(1{+}3\) split appears at \(D=4\)
  --- time extent \(3.5\), three space extents \(9.996\) each with
  \textbf{spatial isotropy \(0.000\)} across 10 seeds --- but the split
  is the signature, not the dynamics, and \emph{time contracts} (extents
  below the space block). At \(D=10\) all nine space dimensions expand
  (no ``3'').
\item
  \textbf{Real-time dynamics (growing clock \(T\)):} the pure-bosonic
  Lorentzian action is unbounded below --- extents run away
  (\(2000\)--\(8700\)) with no equilibrium (drift \(0.13\)--\(0.15\)).
  This mirrors the known fact that Lorentzian IKKT requires the fermion
  determinant as regulator {[}\hyperref[ref-3]{3},
  \hyperref[ref-4]{4}{]}; a CPU toy without it has no controlled
  configuration space.
\item
  \textbf{Bounded (saturating) noncommutativity}
  \(-\sum_{i<j} g\frac{x_{ij}}{1+x_{ij}}\), \(x=\|\mathrm{comm}\|^2\):
  equilibrium exists and is perfectly symmetric (\(\mathrm{top3}\)-gap
  \(1.06\pm0.01\)); a hard-wired 3-of-9 asymmetry is detectable in this
  regime (\(1.14\)) --- so the negative is a \emph{testable} negative.
  No \(k\)-of-9 selection occurs.
\end{enumerate}

\textbf{Conclusion:} no configuration without \(\hat v\) (or a signature
choice) breaks SO(\(D\)) at toy scale; the only ``1+3'' obtainable is
engineered. Consistent with the known physics: Euclidean IKKT does not
dimensionally reduce {[}\hyperref[ref-15]{15}{]}.

\subsubsection{4.6 Causal-set dimensions: percolation negative,
light-cone
invariance}\label{causal-set-dimensions-percolation-negative-light-cone-invariance}

\textbf{Calibration.} The relation-fraction estimator \(\rho \to d\) was
calibrated on known posets per \(N\): chain (\(d{=}1\)), Minkowski
sprinklings \(d{=}2..8\) (\(\rho(2){=}0.25\), \(\rho(4){=}0.086\),
\(\rho(8){=}0.011\)), stable across \(N{=}250,500,1000\).

\textbf{Sequential-growth poset.} The real growth output (relation
\(|\hat v_u\cdot\hat v_k|>0.9\)) is a total order:
\textbf{\(d_{\mathrm{MM}}=1.00\pm0.00\)} at \(M=8,16\). Birth-order-only
time is exactly 1-dimensional; there is no spatial structure in the
current states (\(\mathrm{Tr}((X_u-X_k)^2)\) is nearly constant across
pairs).

\textbf{Random percolation is a finite-size artifact.} Bernoulli(\(p\))
links plus transitive closure at \(N=250\) produced a promising-looking
``\(d\approx4\)'' at \(p=0.016\) (\(d=4.42\)). The scale audit destroys
it:

{\def\LTcaptype{none} % do not increment counter
\begin{longtable}[]{@{}ll@{}}
\toprule\noalign{}
scaling (N=250/500/1000) & d\_MM \\
\midrule\noalign{}
\endhead
\bottomrule\noalign{}
\endlastfoot
fixed \(p=0.02\) & 3.40 / 1.86 / \textbf{1.51} \\
fixed \(k=pN=5\) & 3.40 / 4.48 / \textbf{5.35} \\
\end{longtable}
}

No scaling invariance in either variable (std across \(N\):
\(0.26\)--\(0.82\)). \emph{Mechanism:} the closure of a random DAG at
large \(N\) degenerates to a near-total-order (fixed \(p\)) or has
\(\rho\) scaling differently from the calibration sprinklings (fixed
\(k\)). Per the pre-registered acceptance criteria this is recorded as a
definitive negative for random percolation as a model of 4D causal sets.

\textbf{The relational light-cone law is scale-invariant --- a
structural discovery.} Replacing random links by the matrix-driven rule
(\(|\hat v_u\cdot\hat v_k|>\theta\) and
\(D_{\mathrm{space}}\le(c\Delta t)^2\), with birth-time rescaled to a
fixed cosmic interval) yields causal dimension \textbf{invariant across
\(N\)}:

{\def\LTcaptype{none} % do not increment counter
\begin{longtable}[]{@{}lll@{}}
\toprule\noalign{}
construction & \(d_{\mathrm{MM}}\) (N=250/500/1000) & std across N \\
\midrule\noalign{}
\endhead
\bottomrule\noalign{}
\endlastfoot
percolation (worst case) & 1.5--5.4 & \textbf{0.26--0.82} \\
\textbf{light-cone walk (d\_sp=1)} & 1.17/1.17/1.17 & \textbf{0.002} \\
light-cone walk (d\_sp=3) & 1.21/1.21/1.21 & 0.002 \\
uniform scatter, cone & 1.29/1.29/1.29 & 0.00 \\
\end{longtable}
}

The std across \(N\) drops by two orders of magnitude relative to
percolation. \textbf{Locality (a light-cone) is the condition that makes
causal dimension a stable property} --- this is the positive content of
the route. However, the absolute dimension is a monotone function of the
spatial-diffusion and light-speed parameters (\(\sigma\), \(c\));
c-scans reach \(d=2\) (uniform, \(d_{\mathrm{spatial}}=1\), \(c=0.2\))
and cross \(d=4\) (3-torus walk, \(c\approx0.1\)--\(0.2\)). The
dimension is parameter-controlled --- again engineered, not emergent ---
but now \emph{stably} so.

\textbf{Commutator-compatibility metric on real growth states ---
invariance without engineered coordinates.} Replacing the state
difference by the commutator distance

\[D_{\mathrm{comp}}(u,k)=\frac{1}{d_{\mathrm{sp}}N}\sum_a \big|\mathrm{Tr}\big([X^a_u,X^a_k]^2\big)\big|\]

with the same cone rule, alignment threshold and transitive closure,
uses states generated \emph{directly by the sequential growth model} ---
no torus embedding, no spatial parameter inserted:

{\def\LTcaptype{none} % do not increment counter
\begin{longtable}[]{@{}
  >{\raggedright\arraybackslash}p{(\linewidth - 4\tabcolsep) * \real{0.3333}}
  >{\raggedright\arraybackslash}p{(\linewidth - 4\tabcolsep) * \real{0.3333}}
  >{\raggedright\arraybackslash}p{(\linewidth - 4\tabcolsep) * \real{0.3333}}@{}}
\toprule\noalign{}
\begin{minipage}[b]{\linewidth}\raggedright
construction (real states)
\end{minipage} & \begin{minipage}[b]{\linewidth}\raggedright
\(d_{\mathrm{MM}}\) (N=250/500/1000)
\end{minipage} & \begin{minipage}[b]{\linewidth}\raggedright
std across N
\end{minipage} \\
\midrule\noalign{}
\endhead
\bottomrule\noalign{}
\endlastfoot
commutator, \(f=0.5\) (\(\theta=0.7\)--\(0.9\)) & 3.51 / 3.52 / 3.50 &
\textbf{0.008} \\
commutator, \(f=1.0\) & 1.89 / 1.89 / 1.90 & 0.004 \\
\end{longtable}
}

The relation fraction is scale-invariant (5--8 seeds) and insensitive to
the observer-alignment threshold \(\theta\) (adjacent alignment is
\(1.000\) in the growth cascade, so the threshold is inert). The
underlying states are near-independent, yet the commutator distance is
bimodal (roughly half of the pairs nearly commute) and
separation-independent; combined with the fixed-interval birth-time
rescaling, this yields an \(N\)-independent relation fraction.
\textbf{Epistemic boundary (stated plainly):} scaling invariance --- the
stability of the geometry --- is real; the dimension \emph{value}
remains parameter-steered (\(d\approx3.5\) at \(f=0.5\) arises from our
state-scale self-calibration of the light speed,
\(c_{\mathrm{ref}}=\sqrt{\overline{D}_{\mathrm{comp}}}/5\); \(f=1.0\)
gives \(1.9\)). Spontaneous dimensional selection is absent in every
construction tested in this work.

\subsubsection{4.7 Spectral dimension on flow
networks}\label{spectral-dimension-on-flow-networks}

On a chain (64 nodes), the symmetric random walk gives the clean
\(d_s\approx0.8\)--\(0.9\) plateau (dimension 1); a biased walk (flow)
first raises \(d_s\) (drift) and then collapses it to \(\approx0\) as
the finite chain reaches stationarity (\(P(t)\to\) constant). On a 3-ary
tree: undirected gives a plateau rising toward the Bethe-lattice value
(\(\approx2\)--\(3\) at our depth); \emph{pure outward flow} (transient
walk) gives return probability \(P(t)\to0\) and \textbf{\(d_s\)
diverges} (\(14\to70\to278\to1110\)). Verdict: on a flow network \(d_s\)
measures drift/transience, not geometry; spectral dimension is only
well-defined on symmetric (equilibrium) structures --- unlike the
plateau claims of causal-dynamical-triangulation programs
{[}\hyperref[ref-14]{14}{]}, which evaluate \(d_s\) on equilibrium
geometries. This closes the ``\(4\)D via \(d_s\)'' hope with a mechanism
(and re-derives the earlier null-compatibility of \(d_s\) in causal
sets). \emph{Technical note:} chains/trees are bipartite, so the raw
return probability vanishes on odd \(t\); the lazy walk
(\(T\to\frac12(I+T)\)) is required before any log-derivative estimate (a
documented trap, Appendix \hyperref[appendix-b-trap-catalog]{B}).

\begin{center}\rule{0.5\linewidth}{0.5pt}\end{center}

\subsection{5. Discussion}\label{discussion}

\textbf{Engineered \(\ne\) emergence.} Every positive result here is a
\emph{construction} whose asymmetry is inserted: a special initial
state, an asymmetric birth order, a directed transport term, a
signature, or a diffusion/light-speed parameter. The audit's value is
precisely that it forces this distinction --- eleven metrics that felt
like emergence were shown to be artifacts or null-compatible.

\textbf{What the quest established.} (i) Equilibrium ensembles of this
toy family cannot host an arrow; (ii) process-based constructions can,
deterministically, and their mechanism is understandable (relaxation
capacity; contraction to the past mean; directed transport); (iii) the
``4D'' target is not reachable at toy scale by any equilibrium or random
construction we tested, and the specific reason in each case matches
known physics (no fermion regulator; no metric; no seed concentration);
(iv) \textbf{relational light-cone locality is the discovered necessary
condition for scale-invariant causal dimension} --- a falsifiable,
transferable structural claim; (v) thermalization sufficiency is the
single most dangerous gate (it reversed two conclusions here, and one
earlier v5→v6 correction {[}\hyperref[ref-7]{7}{]}).

\textbf{Limits.} \(N\le8\) matrices, \(\le1000\) events, single-machine
seeds; the light-cone law's dimension parameterization was not optimized
beyond the scans shown; no fermion-like term was ever included; the
``3'' of the physical dimensional reduction remains unexplained by any
toy here (it requires the phase/fermionic machinery of the Lorentzian
programs {[}\hyperref[ref-3]{3}, \hyperref[ref-4]{4}{]}).

\textbf{What an emergent theory would need (candidate directions).} A
regulated (bounded) noncommutativity that \emph{prefers} a subset of
pairs; or a growth law whose spatial structure is dynamical rather than
parameterized; both would be immediately testable by the same six-step
audit, with \(N\to1000\) scale invariance as the acceptance gate.

\begin{center}\rule{0.5\linewidth}{0.5pt}\end{center}

\subsection{6. Conclusion}\label{conclusion}

In the SGOED toy universe: the arrow of time is not in the equilibrium
ensemble --- it enters through the initial condition or through
asymmetric process rules, reproducibly and deterministically (past
hypothesis \(dS/S_0=0.98\); growth inheritance \(1.000\pm0.000\) via
contraction to the past mean; upwind transport decay \(>5000\times\),
\(J=+5.77\pm0.02\), reversal- and null-audited). The 4D question,
pursued across matrix, bounded, dynamical, percolation, and spectral
frameworks, yields a consistent negative --- with one structural
discovery (light-cone locality → scale-invariant causal dimension) and
one methodological deliverable (the six-step audit + trap catalog) that
transfers to any future attempt.

\begin{center}\rule{0.5\linewidth}{0.5pt}\end{center}

\subsection{Appendix A:
Reproducibility}\label{appendix-a-reproducibility}

All scripts run on CPU with fixed seeds; results JSONs saved alongside.
Key parameters:

\begin{itemize}
\tightlist
\item
  v7 unit: \(N\le8\), \(D=2\), coupling \(g_{XY}=0.8\), gates
  \(\lambda=1\), \texttt{max\_extent=10}.
\item
  Past hypothesis: \texttt{step\_past\_hypothesis.py} (\(N=4\), 3
  seeds).
\item
  Sequential growth: \texttt{step\_sequential\_growth.py} (\(M=8..32\),
  \(g_{\mathrm{inter}}=20\), \(n_{\mathrm{therm}}=100\)); mechanism:
  \texttt{step\_growth\_mechanism.py}; past-hypothesis variant:
  \texttt{step\_growth\_past\_hypothesis.py}
  (\(n_{\mathrm{therm}}=120\)).
\item
  Transport: \texttt{step\_langevin\_transport\_tuned.py} +
  \texttt{step\_transport\_robust.py} (8 seeds \(\times\) 500 steps;
  scan \(g_{\mathrm{trans}}\times g_{\mathrm{sink}}\)).
\item
  4D probes: \texttt{sgoed\_matrix\_v15.py},
  \texttt{step\_v15\_dynamical.py}, \texttt{step\_v15\_bounded.py}.
\item
  Causal sets: \texttt{step\_causal\_set\_dmm.py},
  \texttt{step\_causal\_set\_scale\_study.py} (bitset transitive
  closure; per-\(N\) calibration), \texttt{step\_growth\_lightcone.py},
  \texttt{step\_lightcone\_followup.py}.
\item
  Spectral dimension: \texttt{step\_spectral\_dimension\_flow.py} (lazy
  walk).
\item
  Audit toolbox (Supplement A): \texttt{code/audit\_gates.py} --- six
  gates as reusable functions, with a self-test that demonstrates each
  trap class. The negative-evidence slice behind Sections 3/4.1/Appendix
  B ships in \texttt{audit\_evidence/} (v8--v14 cores, audit scripts,
  results) and \texttt{notes/SGOED\_v8..v14\_notes.md}. The
  bistability/hysteresis claim of Section 2.1 is backed by
  \texttt{audit\_evidence/code/AUDIT\_v7\_hysteresis.py} + its result
  JSON.
\item
  \textbf{Self-tuning variants audited:} degree-feedback (SOC-style,
  \texttt{code/step\_growth\_soc.py}) fails the scale-invariance gate
  outright in all 7 configs (std across \(N\) of 0.25--1.09);
  curvature-feedback (\texttt{c\_eff},
  \texttt{code/step\_growth\_ceff.py}) fails outright in 2 of 12 configs
  (std 0.45--0.48), while its 9 remaining non-control runs sit pinned at
  the calibration ceiling (\(d=8.00\), std \(=0\) --- a boundary
  artifact, 5 of them yielding NaN \(c_{\mathrm{eff}}\) statistics ---
  not genuine invariance); only the static-cone control passes genuinely
  (std 0.006). Per-config results:
  \texttt{results/step\_growth\_soc\_results.json},
  \texttt{results/step\_growth\_ceff\_results.json}; full record in the
  living summary (updates 13--14).
\end{itemize}

Seeds: \(\{42,\dots,53\}\) (≥10 for fine grids; ≥6 elsewhere, as
reported). Detailed reproducibility notes and the full results chain
(updates 1--17) in \texttt{notes/SGOED\_TIME\_EMERGENCE\_SUMMARY.md};
project context in \texttt{notes/SGOED\_PROJECT\_SUMMARY.md}.

\subsection{Appendix B: Trap Catalog}\label{appendix-b-trap-catalog}

Every metric that ``looked exciting'' in this project failed exactly one
of the gates below. The catalog is supplied as a structured eight-trap
table below plus a runnable toolbox (\texttt{code/audit\_gates.py},
Supplement A); the full \textbf{12}-trap catalog and an eight-step
auditing workflow ship as Supplement B
(\texttt{SUPPLEMENT\_B\_AUDIT\_HANDBOOK.md}); the v8--v14 scripts and
results backing each row ship in \texttt{audit\_evidence/}.

{\def\LTcaptype{none} % do not increment counter
\begin{longtable}[]{@{}
  >{\raggedright\arraybackslash}p{(\linewidth - 8\tabcolsep) * \real{0.2000}}
  >{\raggedright\arraybackslash}p{(\linewidth - 8\tabcolsep) * \real{0.2000}}
  >{\raggedright\arraybackslash}p{(\linewidth - 8\tabcolsep) * \real{0.2000}}
  >{\raggedright\arraybackslash}p{(\linewidth - 8\tabcolsep) * \real{0.2000}}
  >{\raggedright\arraybackslash}p{(\linewidth - 8\tabcolsep) * \real{0.2000}}@{}}
\toprule\noalign{}
\begin{minipage}[b]{\linewidth}\raggedright
\#
\end{minipage} & \begin{minipage}[b]{\linewidth}\raggedright
trap
\end{minipage} & \begin{minipage}[b]{\linewidth}\raggedright
symptom
\end{minipage} & \begin{minipage}[b]{\linewidth}\raggedright
gate that caught it
\end{minipage} & \begin{minipage}[b]{\linewidth}\raggedright
fix / lesson
\end{minipage} \\
\midrule\noalign{}
\endhead
\bottomrule\noalign{}
\endlastfoot
1 & Histogram artifact & \(R\approx0.5\) from symmetric graphs;
shuffling the values preserves it & 3. Null & never report \(R\) alone;
always compare to shuffle \\
2 & Labeling dependence & sign-count \(D\) ranges \([-98,+28]\) under
node permutations & 4. Labeling & use permutation-invariant metrics \\
3 & Null-compatible dimension & \(d_{\mathrm{MM}}\), \(d_s\) on
symmetrized graphs match random baselines & 3. Null & calibrate any
dimension estimator on known posets first \\
4 & Thermalization insufficiency & two reversals: v5's \texttt{step=2}
sampler; past-hypothesis growth at \(n_{\mathrm{therm}}=60\) & 5.
Thermalization & verify the observable plateaus at the reported
\(n_{\mathrm{therm}}\) \\
5 & Finite-size ``4D'' & \(d\approx4\) at \(N=250\) evaporates at
\(N=1000\) & 2. + 3. Significance/Null & run the scale audit to
\(N\ge1000\) before any dimension claim \\
6 & Bipartite parity & return-probability \(d_s\) vanishes on odd \(t\);
without the lazy walk \(d_s\) oscillates by \(\pm10^3\) & 5.
Thermalization & use the lazy walk \(T\to\frac12(I+T)\) \\
7 & Integer overflow / closure collapse & \texttt{int8} relation-matrix
matmul overflows for \(N>15\); pure-squaring closure (no union)
collapses to zero --- silent garbage & 1. Reproduce (self-test) &
\texttt{int64} arithmetic; \(R\leftarrow R\vee R^2\) \\
8 & Normalization scale & matrix-encoded distances normalized by
\(1/(KN)\) compress the effective light speed, zeroing the dimension
readout at \(c=1\) & 3. Null / calibration & unit-match (aspect
calibrate) to the estimator's convention \\
\end{longtable}
}

The transferable lesson of the whole quest: \textbf{a metric is a
hypothesis about an observable, and each gate above is a way it can be
falsified.}

\subsection{References}\label{references}

{[}1{]} Page, D. N., \& Wootters, W. K. (1983). Evolution without
evolution: Dynamics described by stationary observables. \emph{Physical
Review D}, 27(12), 2885--2892.

{[}2{]} Connes, A., \& Rovelli, C. (1994). Von Neumann algebra
automorphisms and time-thermodynamics relation. \emph{Classical and
Quantum Gravity}, 11(12), 2899--2918.

{[}3{]} Kim, S.-W., Nishimura, J., \& Tsuchiya, A. (2012). Expanding
(3+1)-dimensional universe from a Lorentzian matrix model for
superstring theory in (9+1) dimensions. \emph{Physical Review Letters},
108(1), 011601.

{[}4{]} Kim, S.-W., Nishimura, J., \& Tsuchiya, A. (2012). Late time
behaviors of the expanding universe in the IIB matrix model.
\emph{Journal of High Energy Physics}, 2012(10), 147.

{[}5{]} Rovelli, C. (1996). Relational quantum mechanics.
\emph{International Journal of Theoretical Physics}, 35(8), 1637--1678.

{[}6{]} Chanpengpad, S. (2026). \emph{Structural Time Framework}.
Zenodo.

{[}7{]} Chanpengpad, S. (2026). \emph{SGOED: An Atemporal Matrix
Formalism for Stability-Gated Crystallization and Aggregation (Phase
1--2)}. Zenodo.
\href{https://doi.org/10.5281/zenodo.21786260}{doi:10.5281/zenodo.21786260}

{[}8{]} Metropolis, N., et al.~(1953). Equation of state calculations by
fast computing machines. \emph{Journal of Chemical Physics}, 21(6),
1087--1092.

{[}9{]} Barbour, J. (1999). \emph{The End of Time}. Oxford University
Press.

{[}10{]} Rovelli, C. (2018). \emph{The Order of Time}. Riverhead Books.

{[}11{]} Rideout, D. P., \& Sorkin, R. D. (2000). Classical sequential
growth dynamics for causal sets. \emph{Physical Review D}, 61(2),
024002.

{[}12{]} Sorkin, R. D. (2005). Causal sets: Discrete gravity. In
\emph{Lectures on Quantum Gravity} (pp.~305--327). Springer.

{[}13{]} Myrheim, J. (1978). Statistical geometry. CERN preprint
TH-2538.

{[}14{]} Ambjørn, J., Jurkiewicz, J., \& Loll, R. (2005). Reconstructing
the universe. \emph{Physical Review D}, 72(6), 064014.

{[}15{]} Ishibashi, N., Kawai, H., Kitazawa, Y., \& Tsuchiya, A. (1997).
A large-\(N\) reduced model as superstring. \emph{Nuclear Physics B},
498(1--2), 467--491.

{[}16{]} Albert, D. Z. (2000). \emph{Time and Chance}. Harvard
University Press.
